\section{Conclusions}
\label{sec:conclusion}

The question this benchmark asks is whether a learned method should replace a
well-implemented classical one in the GNSS coordinate-series pipeline. The
answer is task-specific, and the first version's uniform negative verdict does
not survive a protocol that removes three confounds: a truth inside the
parametric basis, a normalisation mismatch at inference, and a downstream noise
model re-estimated on denoised output.

Denoising is conditional on the truth family. The trajectory model and the RTS
smoother remove about four-fifths of the injected error when the truth is a
trend plus constant harmonics, but they add error when it is not. The learned
denoiser is nearly family-invariant, $0.366 \pm 0.002$ in-family and
$0.363 \pm 0.001$ out-of-family, and under per-window inference it beats every
model-free decomposition in both families; it stays well below the parametric
estimator when the truth lies in its basis. The earlier placement of the U-Net
below SSA and the wavelet came from whole-series normalisation at inference, not
from the network.

Gap filling is the task where a learned method gains the most under the
revised protocol, but the gain is confined to synthetic records; on real data
the learned filler is at parity at best. With per-window inference, which
matches the statistics used in
training, the U-Net is the best method at 30-day synthetic gaps ($3.07$ against
$3.34$~mm for linear interpolation). On the full held-out set of 201 component
series it reaches $2.964$~mm on NGL against $2.967$ for linear interpolation, a
tie, and $3.293$ on JPL against $3.233$, $0.06$~mm behind; by region it is
ahead only in Asia and Europe. It is slightly behind SSA at 5 days and behind
the trajectory model at 120 days. A TCN and the smaller BiGRU match or beat it
at 5 and 30 days, so the advantage is not specific to the encoder--decoder
design, though the BiGRU costs two orders of magnitude more to train. The
earlier failure was an inference defect, not an architectural limit; the
corrected result is a competitive learned method at intermediate gaps, not a
demonstrated improvement over interpolation on real data.

Offset detection remains a classical win at the operating point. The learned
field has higher average precision at large offsets, but its validation-selected
$F_1$ is lower at every amplitude, and on real catalogue labels it trails the
classical detector: $0.457 \pm 0.027$ for the classical arm, $0.232 \pm 0.107$
for the twenty U-Net seeds and $0.273 \pm 0.096$ for the ten TCN seeds, with
four of the thirty seeds collapsing. The TCN is ahead of the U-Net, so the
negative detection verdict does not depend on one architecture, but it is still
below the classical detector. The seed spread is part of the result, not noise
to be averaged away.

The velocity product narrows the case for learned pre-processing without
reversing it. Under protocol-matched inference no method improves the velocity:
the untreated series and classical denoisers sit at
$0.0586$--$0.0616$~mm\,yr$^{-1}$, the learned denoiser at $0.0605$ and the
learned imputer at $0.0586$. Calibration is likewise preserved
($0.618$--$0.649$ across methods). The earlier factor-of-$3.6$ degradation and
the calibration collapse to $0.197$ were produced by whole-series
normalisation, and they disappear when the network is scored window by window.
The random-walk arm lowers absolute accuracy for every method without changing
the ordering. The practical recommendation is therefore narrower than before:
neither learned nor classical denoising is a substitute for a joint
trajectory/noise fit upstream of a published velocity, and a learned gap filler
is a defensible alternative at 30-day synthetic gaps but has not shown an
advantage over linear interpolation on the full held-out real set. Learned
pre-processing carries no product-level penalty. The remaining product-level objection is that
it carries no gain either.

Two methodological lessons carry beyond this benchmark. A normalisation
mismatch between training and inference can mimic a capacity failure, so learned
pipelines must be scored under the protocol they were trained with before a
negative result is attributed to the model. And a downstream calibration
comparison is only as meaningful as the noise model attached to each method;
re-estimating that model on a denoised output is not a valid protocol, because
the denoiser has removed the information the estimator needs.

What remains open is a truth observed rather than constructed, a multi-station
treatment of common-mode error, and a learned detector trained and validated on
real catalogue labels with a pre-registered operating point. The benchmark and
its result files make those extensions measurable on the same frozen split.
